\chapterphoto{UCSC Crown/Merrill 公交站旁的林线}{photo-chap3}
\chapter{应用建议}

\section{机械方案与软件方案}
\label{sec:mechvssw}

虽然本书聚焦于控制工程在 FRC 中的应用，但要打造一支成功的队伍，控制并非必需；FRC 主要是一项机械竞赛，而不是软件竞赛。我们花费六个多星期来设计和建造机器人，却常常只给软件测试留出两天时间。尽管如此，各队伍依然能派出有竞争力的机器人——这也证明了如今 FRC 软件生态是多么简单易用。首先应当专注于可靠的机械设计，因为好的机械设计配上简单的软件就能成功，而不可靠的硬件则无论如何都无法让机器人成功。

一个设计问题的解决方案，可能是机械复杂度与软件复杂度之间的权衡。例如，对于一个只需要两个位置的机构，与``电机 + 旋转编码器 + 软件反馈控制''相比，一个连接到气动执行器的电磁阀可能更省事、也更可靠。不过，如果软件方案能够顺利实现，机器人就可以省去空气压缩机和储气罐所占的空间和重量。

我评估设计的经验法则是：优先选择优雅的机械方案，而不是面面俱到的软件方案，因为在比赛中让机械方案变得可靠要更容易。布置合理的传感器也能大幅提升机器人的表现，并减轻操作手的认知负担。例如在 2011 年 FRC 比赛中，用一个限位开关加比赛计时器，就能在终局（endgame）开始时自动释放 minibots。许多情况下，手动流程可以之后再自动化，并且在相关软件或传感器失效时保留手动后备方案。对于必须依赖闭环控制才能工作的设计要保持谨慎。

什么样的问题应当用硬件而不是用巧妙的控制软件来解决？控制可以应对诸如电池电压跌落或测量噪声之类的扰动，但它的能力是有限度的。例如，面对底盘变速箱卡死或者掉链条，软件是无能为力的。有时候，直接在硬件上修好问题的根源才是更好的选择。

设计机器人机构时要考虑可控性（controllability）。FRC 971 队的``Mechanical Design for Controllability''研讨会对这个问题有更详细的论述。从中总结出的两个要点是：

\begin{itemize}
  \item 减小齿轮箱的回差（backlash）
  \item 选择能够提供足够控制能力的电机和减速比
\end{itemize}

\begin{media}[视频资源]
``Spartan Series / Mechanical Design for Controllability''（时长 1 小时 31 分钟）\\
Travis Schuh\\
\url{https://youtu.be/VNfFn-gcfFI}
\end{media}

请记住，``用软件修补''并不总是答案。本书剩余的章节假设你已经完成了工程分析，并得出结论：你所选的设计会从更复杂的控制中获益，而且你有时间或有专长把它实现出来。

\section{执行器饱和}
\label{sec:saturation}

反馈控制器根据参考值与当前状态之间的误差来计算输出。现实世界中的被控对象，其可供反馈控制器使用的控制能力并不是无限的。当达到执行器的限制时，反馈控制器的表现就好像增益被暂时降低了（也就是说，它的控制能力下降了）。

我们用一点数学来解释这一点。设控制器为 $u = k(r - x)$，其中 $u$ 是控制量，$k$ 是增益，$r$ 是参考值，$x$ 是当前状态。令 $u_{\max}$ 为执行器输出的上限（它小于 $u$ 未封顶时的取值），$k_{\max}$ 为与之对应的最大增益。现在针对相同的参考值和当前状态，比较封顶与未封顶的控制器：

\begin{align*}
  u_{\max} &< u \\
  k_{\max} (r - x) &< k (r - x) \\
  k_{\max} &< k
\end{align*}

要使不等式成立，$k_{\max}$ 必须小于原来的 $k$。当系统响应中状态随时间呈线性变化、而不是在接近参考值时呈指数变化，就能明显看出这种增益的降低。这是因为此时控制量不再遵循衰减指数曲线。一旦系统更接近参考值，控制器就会停止饱和，重新给出符合实际的控制器输出。
